\documentclass[a4paper,12pt]{article}

\usepackage{polyglossia}
\setmainlanguage{polish}
\usepackage{geometry}
\geometry{top=2.5cm, bottom=2.5cm, left=3cm, right=3cm}
\usepackage{graphicx}
\usepackage{hyperref}
\usepackage{listings}
\usepackage{xcolor}
\usepackage{float}
\usepackage{caption}
\usepackage{enumitem}
\usepackage{amsmath}
\usepackage{fancyhdr}
\setlength{\headheight}{14.5pt}
\addtolength{\topmargin}{-2.5pt}
\usepackage{titlesec}
\usepackage{fontspec}
\setmainfont{DejaVu Serif}
\setsansfont{DejaVu Sans}
\setmonofont{DejaVu Sans Mono}

\definecolor{codegreen}{rgb}{0,0.6,0}
\definecolor{codegray}{rgb}{0.5,0.5,0.5}
\definecolor{codepurple}{rgb}{0.58,0,0.82}
\definecolor{backcolour}{rgb}{0.95,0.95,0.92}

\lstdefinestyle{csharpstyle}{
    language=[Sharp]C,
    backgroundcolor=\color{backcolour},
    commentstyle=\color{codegreen},
    keywordstyle=\color{magenta},
    numberstyle=\tiny\color{codegray},
    stringstyle=\color{codepurple},
    basicstyle=\ttfamily\small,
    breakatwhitespace=false,
    breaklines=true,
    captionpos=b,
    keepspaces=true,
    numbers=left,
    numbersep=5pt,
    showspaces=false,
    showstringspaces=false,
    showtabs=false,
    tabsize=4,
    frame=single,
    rulecolor=\color{codegray}
}

\lstset{style=csharpstyle}

\hypersetup{
    colorlinks=true,
    linkcolor=blue,
    urlcolor=blue,
    citecolor=blue
}

\pagestyle{fancy}
\fancyhf{}
\fancyhead[L]{VirusSimulation}
\fancyhead[R]{\thepage}
\renewcommand{\headrulewidth}{0.4pt}

\titleformat{\section}{\Large\bfseries}{\thesection.}{1em}{}
\titleformat{\subsection}{\large\bfseries}{\thesubsection.}{1em}{}
\titleformat{\subsubsection}{\normalsize\bfseries}{\thesubsubsection.}{1em}{}

\begin{document}

\begin{titlepage}
    \centering
    \vspace*{4cm}
    {\Huge\bfseries VirusSimulation\par}
    \vspace{0.5cm}
    {\Large Symulacja epidemii w konsoli\par}
    \vspace{0.5cm}
    {\large --- Dokumentacja ---\par}
    \vfill
\end{titlepage}

\tableofcontents
\newpage

\section{Wstęp}

\subsection{Cel projektu}

Celem projektu jest stworzenie konsolowej symulacji rozprzestrzeniania się wirusa na dwuwymiarowej,
proceduralnie generowanej mapie. Symulacja uwzględnia różne typy ludności (pracownicy, studenci,
seniorzy, lekarze), wirusy o zróżnicowanych parametrach oraz mechanizmy zapobiegawcze, takie jak
blokada regionów i świadomość społeczna.

\subsection{Wymagania techniczne}

\begin{itemize}
    \item \textbf{Język:} C\# 13.0
    \item \textbf{Platforma:} .NET 10.0
    \item \textbf{Testy:} xUnit 2.x
    \item \textbf{System kontroli wersji:} Git
    \item \textbf{CI:} GitHub Actions (sprawdzanie formatowania kodu)
\end{itemize}

\subsection{Funkcjonalności}

\begin{itemize}
    \item Proceduralne generowanie mapy z czterema regionami oddzielonymi wodą
    \item Trzy typy wirusów o różnych parametrach: Flu, Covid, Rabies
    \item Pięć typów jednostek: Worker, Student, Elder, Doctor, Hospital
    \item Trzy strategie rozprzestrzeniania się wirusa: jednolita, odległościowa, skupiona
    \item System leczenia (lekarze mobilni i szpitale stacjonarne)
    \item Mechanizmy zapobiegawcze: blokada regionów, świadomość społeczna
    \item Migracja ludności między regionami
    \item Eksport statystyk do pliku CSV
    \item Wizualizacja w konsoli z kolorami ANSI
\end{itemize}

\section{Struktura projektu}

\subsection{Układ katalogów}

\begin{lstlisting}[frame=single, numbers=none, basicstyle=\ttfamily\small]
VirusSimulation/
  Interfaces/       - Interfejsy (ITile, IVirus, ISpreadLogic, IHealingAbility, IRandom)
  Models/
    Entities/       - Entity (bazowa), Hospital
    Humans/         - Human (bazowa), Worker, Student, Elder, Doctor
    Map/            - Tile (enum)
    SpreadLogic/    - Strategie rozprzestrzeniania
    Viruses/        - Flu, Covid, Rabies
  Utilities/        - Board, MapGenerator, Rendering, Statistics, GameRandom
  tests/            - Projekt testowy xUnit
  *.puml            - Diagramy PlantUML (klas, obiektów, sekwencji, stanów)
\end{lstlisting}

\subsection{Pliki konfiguracyjne}

\begin{itemize}
    \item \texttt{VirusSimulation.csproj} --- główny plik projektu (.NET 10.0, exe)
    \item \texttt{tests/VirusSimulation.Tests/VirusSimulation.Tests.csproj} --- projekt testowy
    \item \texttt{tests/VirusSimulation.Tests/xunit.runner.json} --- konfiguracja testów
    \item \texttt{.github/workflows/dotnet-format.yml} --- CI w GitHub Actions
\end{itemize}

\section{Architektura i wzorce projektowe}

Projekt został zaprojektowany z zachowaniem zasad programowania obiektowego i wykorzystuje
kilka kluczowych wzorców projektowych.

\subsection{Hierarchia dziedziczenia}

Rysunek~\ref{fig:class-diagram} przedstawia pełny diagram klas. Główna hierarchia dziedziczenia:

\begin{center}
\begin{lstlisting}[frame=none, numbers=none, basicstyle=\ttfamily\small]
Entity (bazowa klasa dla wszystkich obiektów na planszy)
  +-- Human (bazowa klasa dla ludzi)
  |     +-- Worker    (podwójny ruch)
  |     +-- Student   (50% odporności)
  |     +-- Elder     (50% szansy na ruch)
  |     +-- Doctor    (ruch + leczenie)
  +-- Hospital (stacjonarna jednostka lecznicza)
\end{lstlisting}
\end{center}

\subsection{Zastosowane wzorce projektowe}

\subsubsection{Wzorzec Strategii (Strategy)}

Wirusy nie dziedziczą logiki rozprzestrzeniania --- komponują ją poprzez interfejs
\texttt{ISpreadLogic}. Każdy wirus zawiera instancję konkretnej strategii:

\begin{itemize}
    \item \texttt{Flu} $\rightarrow$ \texttt{UniformSpreadLogic}
    \item \texttt{Covid} $\rightarrow$ \texttt{DistanceWeightedSpreadLogic}
    \item \texttt{Rabies} $\rightarrow$ \texttt{FocusedSpreadLogic}
\end{itemize}

Dzięki temu można dowolnie łączyć wirusy ze strategiami bez modyfikacji istniejących klas.

\subsubsection{Metoda Szablonowa (Template Method)}

Klasa abstrakcyjna \texttt{SpreadLogicBase} definiuje szkielet metody \texttt{Spread()}:

\begin{enumerate}
    \item Iteracja po wszystkich zarażonych ludziach
    \item Dla każdego źródła --- wywołanie abstrakcyjnej metody \texttt{SpreadFromSource()}
    \item Sprawdzenie śmiertelności --- losowanie względem \texttt{virus.Mortality}
\end{enumerate}

Klasy pochodne implementują jedynie \texttt{SpreadFromSource()}, co wymusza spójną strukturę
przy jednoczesnej elastyczności.

\subsubsection{Wzorzec Fabryki (Factory)}

\texttt{SpreadLogicFactory} tworzy i cache'uje singletonowe instancje strategii:

\begin{lstlisting}[language={[Sharp]C}]
public static ISpreadLogic Create(SpreadLogicType type)
{
    return type switch
    {
        SpreadLogicType.Uniform => _uniform ??= new UniformSpreadLogic(),
        SpreadLogicType.DistanceWeighted => _distance ??= new DistanceWeightedSpreadLogic(),
        SpreadLogicType.Focused => _focused ??= new FocusedSpreadLogic(),
        _ => throw new ArgumentException($"Unknown spread logic type: {type}")
    };
}
\end{lstlisting}

\subsubsection{Wstrzykiwanie Zależności (Dependency Injection)}

Interfejs \texttt{IRandom} jest wstrzykiwany do wszystkich klas wymagających losowości
(\texttt{Board}, \texttt{Human}, \texttt{SpreadLogicBase}). W produkcji używany jest
\texttt{GameRandom.Create()} (brak ziarna), w testach \texttt{GameRandom.Create(seed)}
(deterministyczne, powtarzalne wyniki).

\subsubsection{Segregacja Interfejsów (Interface Segregation)}

Zamiast jednego rozbudowanego interfejsu, projekt zawiera pięć wyspecjalizowanych:

\begin{itemize}
    \item \texttt{ITile} --- symbol, opis, priorytet renderowania
    \item \texttt{IVirus} --- właściwości wirusa, metoda rozprzestrzeniania
    \item \texttt{ISpreadLogic} --- strategia rozprzestrzeniania
    \item \texttt{IHealingAbility} --- zdolność leczenia
    \item \texttt{IRandom} --- generator liczb losowych
\end{itemize}

\subsection{Pętlą symulacji}

Główna pętla w \texttt{Program.cs} składa się z dwóch faz na każdy tick:

\begin{enumerate}
    \item \textbf{Aktualizacja jednostek} --- każda jednostka na planszy wywołuje
          \texttt{Update(Board)}. Ludzie poruszają się, lekarze i szpitale leczą.
    \item \textbf{Rozprzestrzenianie wirusa} --- aktywny wirus wywołuje \texttt{Spread(Board)},
          który poprzez strategię rozprzestrzeniania infekuje nowe ofiary i sprawdza śmiertelność.
\end{enumerate}

Warunki zakończenia symulacji:
\begin{itemize}
    \item Liczba zarażonych osiąga 0 (wirus wyeliminowany)
    \item Wszyscy ludzie są zarażeni
    \item Liczba zarażonych nie zmienia się przez 100 kolejnych ticków (stagnacja)
\end{itemize}

\section{Opis komponentów}

\subsection{Interfejsy (\texttt{Interfaces/})}

\begin{itemize}
    \item \texttt{ITile} --- definiuje \texttt{Symbol} (znak), \texttt{Description} (opis),
          \texttt{RenderPriority} (priorytet renderowania). Implementowany przez \texttt{Entity}.
    \item \texttt{IVirus} --- definiuje \texttt{Name}, \texttt{Infectivity}, \texttt{Mortality},
          \texttt{InfectionRange} oraz metodę \texttt{Spread(Board)}.
    \item \texttt{ISpreadLogic} --- kontrakt dla strategii rozprzestrzeniania: \texttt{Spread(IVirus, Board)}.
    \item \texttt{IHealingAbility} --- definiuje \texttt{Range}, \texttt{Effectiveness},
          \texttt{Heal(Human)}. Implementowany przez \texttt{Doctor} i \texttt{Hospital}.
    \item \texttt{IRandom} --- abstrakcja generatora liczb losowych: \texttt{Next()},
          \texttt{Next(int)}, \texttt{Next(int, int)}, \texttt{NextSingle()}.
\end{itemize}

\subsection{Jednostki (\texttt{Models/Entities/})}

\subsubsection{Entity}

Klasa bazowa dla wszystkich obiektów na planszy. Przechowuje współrzędne $(x, y)$,
definiuje wirtualne właściwości \texttt{Symbol}, \texttt{Description}, \texttt{RenderPriority}
oraz wirtualną metodę \texttt{Update(Board)}.

\subsubsection{Hospital}

Stacjonarna jednostka lecznicza (symbol \texttt{+}). Implementuje \texttt{IHealingAbility}:
\begin{itemize}
    \item Zasięg: 2
    \item Skuteczność: 1.0
    \item Mechanizm: w każdej aktualizacji znajduje zarażonych w zasięgu i zwiększa
          ich licznik \texttt{HealingTicks}; po 5 tickach następuje pełne wyleczenie.
\end{itemize}

\subsection{Ludzie (\texttt{Models/Humans/})}

\subsubsection{Human}

Klasa bazowa dla wszystkich typów ludzi. Kluczowe pola i metody:

\begin{itemize}
    \item \texttt{IsAlive} --- czy żyje
    \item \texttt{Virus} --- wirus, którym jest zarażony (lub \texttt{null})
    \item \texttt{IsInfected} --- czy jest zarażony
    \item \texttt{HealingTicks} --- licznik postępu leczenia
    \item \texttt{Age} --- wiek (wpływa na migrację)
    \item \texttt{HomeRegion} --- region macierzysty
    \item \texttt{MigrationCooldown} --- czas do kolejnej migracji
\end{itemize}

\textbf{Mechanika ruchu:}
\begin{enumerate}
    \item Ruch lokalny: o 1 krok w losowym kierunku $(\Delta x, \Delta y \in \{-1,0,1\})$
    \item Redukcja przez świadomość: ruch tylko gdy \texttt{rng.NextSingle() > AwarenessLevel}
    \item Blokada: podczas lockdownu nie można opuścić swojego regionu
    \item Migracja dalekodystansowa: losowa teleportacja do innego regionu z cooldownem 15 ticków
\end{enumerate}

Prawdopodobieństwo migracji:
\[
P_{\text{migracji}} = 0.03 \cdot M_{\text{zawodu}} \cdot \left(1.4 - \frac{\text{Age}}{100}\right)
\]

gdzie $M_{\text{zawodu}}$ wynosi: Student 1.4, Worker 1.1, Doctor 1.0, Elder 0.7.
Mnożnik wieku jest clampowany do przedziału $[0.6, 1.4]$.

\subsubsection{Worker}

\begin{itemize}
    \item Symbol: \texttt{W}
    \item Wiek: 23--61
    \item Specjalizacja: wykonuje \textbf{dwa ruchy} na tick (dwukrotne wywołanie \texttt{Move()})
\end{itemize}

\subsubsection{Student}

\begin{itemize}
    \item Symbol: \texttt{S}
    \item Wiek: 16--26
    \item Specjalizacja: \textbf{50\% naturalnej odporności} --- metoda \texttt{Infect()} ma 50\% szansy
          na odrzucenie infekcji
\end{itemize}

\subsubsection{Elder}

\begin{itemize}
    \item Symbol: \texttt{O}
    \item Wiek: 65--91
    \item Specjalizacja: działa tylko w \textbf{50\% aktualizacji} (pomija swoją turę gdy
          \texttt{rng.NextSingle() < 0.5f})
\end{itemize}

\subsubsection{Doctor}

\begin{itemize}
    \item Symbol: \texttt{D}
    \item Implementuje \texttt{IHealingAbility}: zasięg 3, skuteczność 0.8
    \item W każdej aktualizacji: najpierw się przemieszcza, następnie leczy zarażonych w zasięgu
    \item Leczenie trwa 5 ticków (addytywne z innymi źródłami leczenia)
\end{itemize}

\subsection{Wirusy (\texttt{Models/Viruses/})}

\begin{table}[H]
\centering
\begin{tabular}{|l|c|c|c|c|}
\hline
\textbf{Wirus} & \textbf{Infectivity} & \textbf{Mortality} & \textbf{Zasięg} & \textbf{Strategia} \\ \hline
Flu    & 0.25 & 0.005 (0.5\%)  & 1 & Uniform \\
Covid  & 0.65 & 0.02 (2\%)     & 2 & DistanceWeighted \\
Rabies & 0.50 & 0.30 (30\%)    & 1 & Focused \\ \hline
\end{tabular}
\caption{Parametry wirusów}
\end{table}

Każdy wirus deleguje metodę \texttt{Spread()} do swojej strategii rozprzestrzeniania
(\textbf{wzorzec Strategii}).

\subsection{Strategie rozprzestrzeniania (\texttt{Models/SpreadLogic/})}

\subsubsection{SpreadLogicBase}

Klasa abstrakcyjna implementująca \textbf{metodę szablonową}:

\begin{lstlisting}[language={[Sharp]C}]
public void Spread(IVirus virus, Board board)
{
    var infected = board.GetInfectedHumans();
    foreach (var source in infected)
    {
        SpreadFromSource(virus, board, source);

        // Sprawdzenie smiertelnosci
        if (rng.NextSingle() <= virus.Mortality)
            source.Die();
    }
}

protected abstract void SpreadFromSource(IVirus virus, Board board, Human source);
\end{lstlisting}

\subsubsection{UniformSpreadLogic}

Każdy zdrowy człowiek w zasięgu ma taką samą szansę na zarażenie:
\[
P_{\text{infekcji}} = \text{Infectivity}
\]

\subsubsection{DistanceWeightedSpreadLogic}

Szansa na zarażenie maleje liniowo z odległością:
\[
P_{\text{infekcji}} = \text{Infectivity} \cdot \left(1 - \frac{\text{distance}}{\text{range} + 1}\right)
\]

gdzie \texttt{distance} to odległość euklidesowa (zaokrąglona w górę).

\subsubsection{FocusedSpreadLogic}

Wybierany jest jeden losowy cel w zasięgu, który otrzymuje szansę na zarażenie równą
\texttt{Infectivity}.

\subsection{Plansza i generowanie mapy}

\subsubsection{Board}

Centralny stan gry. Zawiera:

\begin{itemize}
    \item Dwuwymiarową tablicę kafelków \texttt{Tile[,] grid} o rozmiarze $40 \times 40$
    \item Listę jednostek \texttt{List<Entity> entities}
    \item Właściwości: \texttt{AwarenessLevel}, \texttt{LockdownEnabled}
\end{itemize}

Kluczowe metody:

\begin{itemize}
    \item \texttt{IsWalkable(int x, int y)} --- sprawdza czy kafelek jest przechodni
    \item \texttt{GetHumansInRange(int cx, int cy, int range)} --- znajdowanie ludzi w zasięgu
          kołowym (wykorzystuje kwadrat odległości euklidesowej)
    \item \texttt{GetInfectedHumans()} --- zwraca listę zarażonych, żywych ludzi
    \item \texttt{addRandomEntity(Func<int,int,Entity> factory)} --- umieszcza jednostkę
          na losowym kafelku lądowym
\end{itemize}

\subsubsection{MapGenerator}

Proceduralne generowanie mapy z wykorzystaniem algorytmu \textbf{Voronoia z BFS}:

\begin{enumerate}
    \item Wybór 4 punktów zarodkowych (seed points) z minimalnym odstępem
          $\text{regionRadius} \cdot 2 \cdot 0.5$
    \item Ekspansja BFS z użyciem \texttt{PriorityQueue} (priorytetem jest odległość)
    \item Strefa buforowa wokół granic regionów zapobiega ich bezpośredniemu sąsiedztwu
    \item Nieprzypisane komórki stają się wodą (\texttt{Water})
    \item Każdy region otrzymuje unikalną wartość enum (\texttt{Region0}--\texttt{Region3})
\end{enumerate}

\subsection{Renderowanie}

Klasa \texttt{Rendering} odpowiada za wyświetlanie stanu symulacji w konsoli z użyciem
kolorów ANSI:

\begin{itemize}
    \item Zarażeni ludzie --- kolor czerwony
    \item Martwi ludzie --- kolor czarny
    \item Lekarze --- jasny zielony
    \item Szpitale --- żółty
    \item Woda --- jasny niebieski
    \item Ląd --- domyślny
\end{itemize}

Każda komórka renderuje symbol jednostki (jeśli występuje) lub symbol kafelka.

\subsection{Statystyki}

Klasa \texttt{Statistics} kolekcjonuje dane w postaci krotek
\texttt{(Tick, Infected, Dead)} i eksportuje je do pliku \texttt{simulation.csv}
po zakończeniu symulacji.

\section{Mechanizmy symulacyjne}

\subsection{Świadomość społeczna (Awareness)}

W miarę wzrostu liczby zarażonych, ludzie zmniejszają swoją mobilność:

\[
\text{AwarenessLevel} = \min\left(1, \frac{\text{liczba zarażonych}}{300}\right)
\]

Ruch następuje tylko gdy \texttt{rng.NextSingle() > AwarenessLevel}. Przy 300 lub więcej
zarażonych świadomość osiąga maksimum (1.0), całkowicie zatrzymując dobrowolny ruch.

\subsection{Blokada (Lockdown)}

Aktywuje się gdy liczba zarażonych przekracza 80\% żyjącej populacji. Podczas lockdownu
ludzie nie mogą przekraczać granic swojego regionu macierzystego.

\subsection{System leczenia}

Zarówno lekarze, jak i szpitale implementują \texttt{IHealingAbility}:

\begin{itemize}
    \item Każde źródło leczenia zwiększa \texttt{HealingTicks} zarażonego o 1 na tick
    \item Po osiągnięciu 5 ticków następuje całkowite wyleczenie (\texttt{Heal()})
    \item Wiele źródeł leczenia działa addytywnie --- np. lekarz i szpital razem
          leczą w $2.5$ ticka
\end{itemize}

\begin{table}[H]
\centering
\begin{tabular}{|l|c|c|}
\hline
\textbf{Jednostka} & \textbf{Zasięg} & \textbf{Skuteczność} \\ \hline
Doctor  & 3 & 0.8 \\
Hospital & 2 & 1.0 \\ \hline
\end{tabular}
\caption{Parametry jednostek leczących}
\end{table}

\section{Testy}

Projekt zawiera zestaw testów jednostkowych w technologii xUnit
(\texttt{tests/VirusSimulation.Tests/}).

\subsection{Konfiguracja}

Testy wykorzystują \texttt{GameRandom.Create(42)} (deterministyczne ziarno) dla
zapewnienia powtarzalności. Równoległe wykonywanie testów jest wyłączone
(\texttt{xunit.runner.json}).

\subsection{Pokrycie testów}

\begin{itemize}
    \item \textbf{HumanTests:} infekcja, leczenie, śmierć, odporność Studentów (test statystyczny
          na 2000 prób), podwójny ruch Workera, 50\% ruchu Eldera (1000 prób)
    \item \textbf{BoardTests:} przechodniość kafelków, pobieranie jednostek, zakresy,
          blokada regionów, losowe rozmieszczanie
    \item \textbf{EntityTests:} leczenie przez szpital (stopniowe i pełne), leczenie przez
          lekarza, właściwości jednostek, \texttt{SpreadLogicFactory}
    \item \textbf{VirusTests:} właściwości wirusów, rozprzestrzenianie na pustej planszy,
          wszystkie trzy strategie (ziarno 42)
    \item \textbf{StatisticsTests:} format CSV, wiele wierszy, wiersze zerowe
\end{itemize}

\subsection{Uruchamianie}

\begin{lstlisting}[frame=single, numbers=none, basicstyle=\ttfamily\small]
dotnet test tests/VirusSimulation.Tests
\end{lstlisting}

\section{Instrukcja uruchomienia}

\subsection{Wymagania}

\begin{itemize}
    \item .NET SDK 10.0 lub nowszy
\end{itemize}

\subsection{Budowanie}

\begin{lstlisting}[frame=single, numbers=none, basicstyle=\ttfamily\small]
dotnet build
\end{lstlisting}

\subsection{Uruchamianie}

\begin{lstlisting}[frame=single, numbers=none, basicstyle=\ttfamily\small]
dotnet run
\end{lstlisting}

Program pyta o:
\begin{enumerate}
    \item Typ wirusa (Flu / Covid / Rabies)
    \item Liczbę poszczególnych jednostek (Workers, Students, Elders, Doctors, Hospitals)
\end{enumerate}

\subsection{Przebieg symulacji}

\begin{enumerate}
    \item Generowanie mapy ($40 \times 40$)
    \item Rozmieszczenie jednostek na lądzie
    \item Zarażenie pierwszej napotkanej osoby
    \item Pętla ticków (z opóźnieniem 200 ms):
          \begin{itemize}
              \item Aktualizacja jednostek
              \item Rozprzestrzenianie wirusa
              \item Renderowanie stanu
              \item Zapis statystyk
          \end{itemize}
    \item Zakończenie po spełnieniu warunku brzegowego
    \item Eksport wyników do \texttt{simulation.csv}
\end{enumerate}

\section{Podsumowanie}

Projekt \texttt{VirusSimulation} stanowi kompleksową implementację symulacji epidemii
w języku C\# z wykorzystaniem zaawansowanych mechanizmów programowania obiektowego:

\begin{itemize}
    \item \textbf{Czytelna architektura} --- hierarchia dziedziczenia, pięć wyspecjalizowanych
          interfejsów
    \item \textbf{Wzorce projektowe} --- Strategia, Metoda Szablonowa, Fabryka, Wstrzykiwanie
          Zależności, Segregacja Interfejsów
    \item \textbf{Testowalność} --- deterministyczne ziarno RNG, interfejs \texttt{IRandom},
          wysoka izolacja komponentów
    \item \textbf{Rozszerzalność} --- łatwość dodawania nowych typów wirusów, strategii
          rozprzestrzeniania, jednostek i mechanik
    \item \textbf{Proceduralna generacja} --- algorytm Voronoia z BFS do tworzenia
          zróżnicowanych map
    \item \textbf{Złożone zachowania} --- migracja, świadomość, blokada, system leczenia,
          odporność, statystyki
\end{itemize}

\appendix

\section{Diagramy UML}

\subsection{Diagram klas}

\begin{figure}[H]
\centering
\includegraphics[width=\textwidth]{class_diagram.png}
\caption{Diagram klas --- hierarchia dziedziczenia, realizacje interfejsów, kompozycje}
\label{fig:class-diagram}
\end{figure}

\subsection{Diagram obiektów}

\begin{figure}[H]
\centering
\includegraphics[width=0.7\textwidth]{object_diagram.png}
\caption{Diagram obiektów --- przykładowy stan symulacji}
\label{fig:object-diagram}
\end{figure}

\subsection{Diagram sekwencji}

\begin{figure}[H]
\centering
\includegraphics[width=\textwidth]{sequence_diagram.png}
\caption{Diagram sekwencji --- dwufazowa pętla symulacji}
\label{fig:sequence-diagram}
\end{figure}

\subsection{Diagram stanów}

\begin{figure}[H]
\centering
\includegraphics[width=0.6\textwidth]{state_diagram.png}
\caption{Diagram stanów człowieka --- cykl życia}
\label{fig:state-diagram}
\end{figure}

\end{document}
